\chapter{§2.10 一次函数}
\label{sec:2.10}


\prerequisite{§2.6 一元一次方程, §2.9 不等式}

\gradelevel{8 年级}


\section*{函数的概念}


\begin{explore}


弹簧实验中，往弹簧下方挂不同质量的重物，测量弹簧的总长度。挂 $1$ 千克时长 $12$ 厘米，挂 $2$ 千克时长 $14$ 厘米，挂 $3$ 千克时长 $16$ 厘米。可以发现：每增加 $1$ 千克，弹簧增长 $2$ 厘米。弹簧长度 $y$ 随重物质量 $x$ 的变化而变化，而且\textbf{每一个 $x$ 的值对应唯一一个 $y$ 的值}——这就是\textbf{函数}关系。

\end{explore}

\begin{understand}


\textbf{变量}：在一个变化过程中，可以取不同值的量叫做变量。

\textbf{常量}：在一个变化过程中保持不变的量。

\textbf{函数}：一般地，在某个变化过程中，如果有两个变量 $x$ 和 $y$ ，对于 $x$ 的每一个取值， $y$ 都有\textbf{唯一确定}的值与之对应，那么我们就说 $y$ 是 $x$ 的函数。

\begin{itemize}
  \item $x$ 叫做\textbf{自变量}。
  \item $y$ 叫做\textbf{因变量}（或函数值）。
  \item 自变量的取值范围叫做函数的\textbf{定义域}。
\end{itemize}


\textbf{为什么强调"唯一确定"？} 这是函数与普通关系的区别。比如 $y^2 = x$ 不是函数关系（一个 $x$ 可能对应两个 $y$ ），但 $y = x^2$ 是函数（一个 $x$ 对应唯一的 $y$ ）。

\end{understand}

\begin{workedexamples}


\textbf{例 1.} 判断下列关系中， $y$ 是否为 $x$ 的函数：

(1) $y = 2x + 1$ 　(2) $y^2 = x$ （ $x \geq 0$ ）　(3) $y = \frac{1}{x}$

\textbf{解：}

(1) 是。每个 $x$ 对应唯一的 $y = 2x + 1$ 。

(2) 不是。当 $x = 4$ 时， $y = 2$ 或 $y = -2$ ，不唯一。

(3) 是。当 $x \neq 0$ 时，每个 $x$ 对应唯一的 $y = \frac{1}{x}$ 。定义域为 $x \neq 0$ 。

\textbf{例 2.} 正方形的面积 $S$ 与边长 $a$ 的关系为 $S = a^2$ 。写出自变量、因变量和定义域。

\textbf{解：} 自变量为 $a$ ，因变量为 $S$ 。因为边长必须为正数，所以定义域为 $a > 0$ 。

\textbf{例 3.} 弹簧实验中， $y = 2x + 10$ （ $x \geq 0$ ），其中 $x$ 为重物质量（千克）， $y$ 为弹簧长度（厘米）。当 $x = 5$ 时， $y$ 等于多少？

\textbf{解：} $y = 2 \times 5 + 10 = 20$ （厘米）。

\end{workedexamples}

\begin{keytakeaway}


\begin{itemize}
  \item 函数的核心：一个 $x$ 对应唯一的 $y$ 。
  \item 函数由三要素确定：解析式（对应法则）、定义域、值域。
  \item 定义域要考虑实际意义和数学限制（分母不为零、被开方数非负等）。
\end{itemize}


\end{keytakeaway}

\begin{practice}


\begin{enumerate}
  \item 圆的面积 $S$ 和半径 $r$ 的关系 $S = \pi r^2$ 中，自变量和因变量分别是什么？定义域呢？
  \item $y = \frac{2}{x - 1}$ 的定义域是什么？
  \item 判断 $\lvert y \rvert = x$ （ $x \geq 0$ ）中 $y$ 是否为 $x$ 的函数。
  \item 若 $y = 3x - 2$ ，求当 $x = -1, 0, 1, 2$ 时对应的 $y$ 值。
  \item 某出租车的收费标准：起步价 $10$ 元（含 $3$ 千米），超出部分每千米 $2$ 元。写出行驶 $x$ 千米（ $x > 3$ ）时的费用 $y$ 与 $x$ 的关系式。
\end{enumerate}


\end{practice}

\section*{函数的表示方法}


\begin{explore}


同一个函数可以用不同的方式来描述。比如弹簧长度与重物质量的关系，可以写成公式 $y = 2x + 10$ ，也可以列一张表格，还可以画一条曲线。这三种方式各有优缺点。

\end{explore}

\begin{understand}


函数的三种表示方法：

\begin{enumerate}
  \item \textbf{解析式法}：用数学公式表示。如 $y = 2x + 1$ 。
\end{enumerate}

   - 优点：精确、简洁、便于计算。
   - 缺点：有些实际关系难以写出公式。

\begin{enumerate}
  \item \textbf{列表法}：用表格列出对应的 $x$ 和 $y$ 值。
\end{enumerate}

   - 优点：直观，查找方便。
   - 缺点：不能列出所有取值。

\begin{enumerate}
  \item \textbf{图象法}：在坐标平面上画出函数的图象。
\end{enumerate}

   - 优点：直观展示变化趋势。
   - 缺点：读数不够精确。

\end{understand}

\begin{workedexamples}


\textbf{例 1.} 函数 $y = x + 2$ ，列出 $x = -2, -1, 0, 1, 2$ 时的函数值表。

\textbf{解：}

\begin{center}
\begin{tabular}{llllll}
\toprule
$x$ & $-2$ & $-1$ & $0$ & $1$ & $2$ \\
\midrule
$y$ & $0$ & $1$ & $2$ & $3$ & $4$ \\
\bottomrule
\end{tabular}
\end{center}


\textbf{例 2.} 根据上表画出 $y = x + 2$ 的图象。

\textbf{解：} 在直角坐标系中描出点 $(-2, 0)$ 、 $(-1, 1)$ 、 $(0, 2)$ 、 $(1, 3)$ 、 $(2, 4)$ ，连成一条直线。这条直线就是 $y = x + 2$ 的图象。

\textbf{例 3.} 某城市某天的气温变化如下表所示。用列表法和图象法分别表示。

\begin{center}
\begin{tabular}{lllllll}
\toprule
时刻 & 6:00 & 9:00 & 12:00 & 15:00 & 18:00 & 21:00 \\
\midrule
气温( $^\circ$ C) & 5 & 10 & 16 & 18 & 12 & 7 \\
\bottomrule
\end{tabular}
\end{center}


\textbf{解：} 列表法已经给出。图象法：以时刻为横轴、气温为纵轴，描出 $6$ 个点，用折线连接，可以直观看出气温先升后降的趋势。

\end{workedexamples}

\begin{keytakeaway}


\begin{itemize}
  \item 函数有三种等价表示方法：解析式、列表、图象。
  \item 解析式最精确，图象最直观，列表最便于查找。
  \item 画函数图象：列表 → 描点 → 连线。
\end{itemize}


\end{keytakeaway}

\begin{practice}


\begin{enumerate}
  \item 函数 $y = -x + 3$ ，列出 $x = -1, 0, 1, 2, 3, 4$ 时的函数值表。
  \item 根据上表描点并画出图象。图象是什么形状？
  \item 函数 $y = x^2$ ，列出 $x = -2, -1, 0, 1, 2$ 时的函数值表，描点画图。图象是直线吗？
  \item 说出解析式法、列表法、图象法各自的一个优点。
  \item 某商品降价促销，第 $1$ 天价格 $100$ 元，以后每天降价 $5$ 元，写出价格 $y$ 与天数 $x$ 的关系式，并画出前 $5$ 天的图象。
\end{enumerate}


\end{practice}

\section*{正比例函数}


\begin{explore}


汽车匀速行驶，每小时走 $60$ 千米。走 $1$ 小时是 $60$ 千米，走 $2$ 小时是 $120$ 千米，走 $t$ 小时是 $60t$ 千米。路程 $s$ 与时间 $t$ 的关系是 $s = 60t$ ——路程与时间成\textbf{正比例}。

\end{explore}

\begin{understand}


\textbf{正比例函数}：形如

\[
y = kx \quad (k \neq 0)
\]


的函数叫做正比例函数。 $k$ 叫做\textbf{比例系数}。

\textbf{正比例函数的图象}：是一条过原点 $(0, 0)$ 的\textbf{直线}。

\begin{itemize}
  \item $k > 0$ ：直线从左下到右上（ $y$ 随 $x$ 增大而增大）。
  \item $k < 0$ ：直线从左上到右下（ $y$ 随 $x$ 增大而减小）。
\end{itemize}


\textbf{为什么 $k \neq 0$ ？} 若 $k = 0$ ， $y = 0$ 是常数函数，不是正比例函数。

\end{understand}

\begin{workedexamples}


\textbf{例 1.} 判断下列函数是否为正比例函数，若是则指出 $k$ ：

(1) $y = 3x$ 　(2) $y = 2x + 1$ 　(3) $y = -\frac{x}{4}$

\textbf{解：}

(1) 是正比例函数， $k = 3$ 。

(2) 不是正比例函数（多了常数项 $+1$ ）。

(3) 是正比例函数， $k = -\frac{1}{4}$ 。

\textbf{例 2.} 画出 $y = 2x$ 和 $y = -x$ 的图象。

\textbf{解：} 正比例函数的图象过原点，只需再找一个点即可确定直线。

\begin{itemize}
  \item $y = 2x$ ：取 $x = 1$ 得 $y = 2$ 。过 $(0,0)$ 和 $(1,2)$ 画直线。
  \item $y = -x$ ：取 $x = 1$ 得 $y = -1$ 。过 $(0,0)$ 和 $(1,-1)$ 画直线。
\end{itemize}


\textbf{例 3.} 已知 $y$ 与 $x$ 成正比例，且当 $x = 3$ 时 $y = 9$ ，求函数解析式。

\textbf{解：} 设 $y = kx$ 。代入 $x = 3, y = 9$ ： $9 = 3k$ ， $k = 3$ 。

所以 $y = 3x$ 。

\end{workedexamples}

\begin{keytakeaway}


\begin{itemize}
  \item 正比例函数 $y = kx$ 的图象是过原点的直线。
  \item $k$ 决定了直线的倾斜方向和程度。
  \item 画正比例函数图象只需两个点：原点和另一个点。
\end{itemize}


\end{keytakeaway}

\begin{practice}


\begin{enumerate}
  \item 写出一个 $k = -2$ 的正比例函数。
  \item $y = 5x$ 的图象经过哪个象限？
  \item 已知 $y$ 与 $x$ 成正比例，且 $x = 2$ 时 $y = -6$ ，求 $y$ 与 $x$ 的关系式。
  \item 正比例函数 $y = kx$ 的图象过点 $(3, -12)$ ，求 $k$ 。
  \item 比较 $y = 3x$ 和 $y = \frac{1}{2}x$ 两条直线，哪条更"陡"？
\end{enumerate}


\end{practice}

\section*{一次函数}


\begin{explore}


出租车的费用与行驶里程的关系为 $y = 2x + 4$ （ $x$ 为超出起步里程的部分， $y$ 为总费用），这不是正比例函数（多了个 $+4$ ），但仍然是一条直线。这种更一般的线性关系就是\textbf{一次函数}。

\end{explore}

\begin{understand}


\textbf{一次函数}：形如

\[
y = kx + b \quad (k \neq 0)
\]


的函数叫做一次函数。 $k$ 是斜率（决定倾斜程度）， $b$ 是截距（图象与 $y$ 轴的交点）。

当 $b = 0$ 时，一次函数就是正比例函数。所以正比例函数是一次函数的特殊情况。

\textbf{一次函数的图象}：是一条\textbf{直线}，与 $y$ 轴的交点为 $(0, b)$ 。

\end{understand}

\begin{workedexamples}


\textbf{例 1.} 画出 $y = 2x + 1$ 的图象。

\textbf{解：} 取两个点：

\begin{itemize}
  \item $x = 0$ 时 $y = 1$ ，得点 $(0, 1)$ 。
  \item $x = 1$ 时 $y = 3$ ，得点 $(1, 3)$ 。
\end{itemize}


过这两点画直线。

\textbf{例 2.} 求一次函数 $y = -3x + 6$ 与 $x$ 轴、 $y$ 轴的交点。

\textbf{解：}

\begin{itemize}
  \item 与 $y$ 轴的交点（令 $x = 0$ ）： $y = 6$ ，交点为 $(0, 6)$ 。
  \item 与 $x$ 轴的交点（令 $y = 0$ ）： $0 = -3x + 6$ ， $x = 2$ ，交点为 $(2, 0)$ 。
\end{itemize}


\textbf{例 3.} 已知一次函数的图象经过 $(1, 5)$ 和 $(3, 11)$ ，求解析式。

\textbf{解：} 设 $y = kx + b$ 。代入两点：

\[
\begin{cases} k + b = 5 \\ 3k + b = 11 \end{cases}
\]


② $-$ ①： $2k = 6$ ， $k = 3$ 。 $b = 5 - 3 = 2$ 。

$y = 3x + 2$ 。

\textbf{例 4.} 一次函数 $y = kx + b$ 中，若 $k > 0$ ， $b < 0$ ，图象经过哪些象限？

\textbf{解：} $k > 0$ 表示直线从左下到右上； $b < 0$ 表示与 $y$ 轴的交点在 $x$ 轴下方。图象经过第一、三、四象限。

\end{workedexamples}

\begin{keytakeaway}


\begin{itemize}
  \item 一次函数 $y = kx + b$ 的图象是直线。
  \item $k$ 决定斜率（倾斜程度和方向）， $b$ 决定与 $y$ 轴的交点。
  \item 确定一条直线只需两个点。
\end{itemize}


\end{keytakeaway}

\begin{practice}


\begin{enumerate}
  \item 画出 $y = -x + 4$ 的图象。
  \item 求 $y = \frac{1}{2}x - 3$ 与 $x$ 轴和 $y$ 轴的交点。
  \item 已知直线经过 $(0, -1)$ 和 $(2, 3)$ ，求一次函数的解析式。
  \item 一次函数 $y = kx + b$ 中 $k < 0$ ， $b > 0$ ，图象经过哪些象限？
  \item 某蜡烛初始长度 $20$ 厘米，每小时燃烧 $2$ 厘米。写出剩余长度 $y$ 与时间 $x$ 的关系式，并指出定义域。
  \item 求两条直线 $y = 2x + 1$ 和 $y = -x + 7$ 的交点。
\end{enumerate}


\end{practice}

\section*{一次函数的图象与性质}


\begin{explore}


在同一个坐标系中画出 $y = 2x$ 、 $y = 2x + 3$ 、 $y = 2x - 1$ ，你会发现它们是三条\textbf{平行}的直线。再画 $y = x + 1$ 和 $y = 3x + 1$ ，它们过同一个点 $(0, 1)$ 但"倾斜程度"不同。 $k$ 和 $b$ 到底如何影响图象？

\end{explore}

\begin{understand}


\textbf{ $k$ 的作用（斜率）}：

\begin{itemize}
  \item $k > 0$ ： $y$ 随 $x$ 的增大而增大（图象"上升"）。
  \item $k < 0$ ： $y$ 随 $x$ 的增大而减小（图象"下降"）。
  \item $\lvert k \rvert$ 越大，直线越陡。
\end{itemize}


\textbf{ $b$ 的作用（截距）}：

\begin{itemize}
  \item $b > 0$ ：图象与 $y$ 轴交于正半轴。
  \item $b = 0$ ：图象过原点（正比例函数）。
  \item $b < 0$ ：图象与 $y$ 轴交于负半轴。
\end{itemize}


\textbf{平行与相交}：

\begin{itemize}
  \item $k$ 相同、 $b$ 不同的两条直线\textbf{平行}。
  \item $k$ 不同的两条直线必\textbf{相交}。
\end{itemize}


\end{understand}

\begin{workedexamples}


\textbf{例 1.} 不画图，说出 $y = -2x + 5$ 的性质。

\textbf{解：}

\begin{itemize}
  \item $k = -2 < 0$ ： $y$ 随 $x$ 增大而减小。
  \item $b = 5 > 0$ ：图象与 $y$ 轴交于 $(0, 5)$ 。
  \item 图象经过第一、二、四象限。
\end{itemize}


\textbf{例 2.} 判断直线 $y = 3x + 1$ 与 $y = 3x - 2$ 的位置关系。

\textbf{解：} $k$ 相同（都是 $3$ ）， $b$ 不同。两条直线\textbf{平行}。

\textbf{例 3.} 若一次函数 $y = kx + b$ 的图象经过第二、三、四象限，求 $k$ 、 $b$ 的符号。

\textbf{解：} 经过第二、三象限说明图象"下降"， $k < 0$ 。经过第三、四象限（不经过第一象限）说明与 $y$ 轴交于负半轴， $b < 0$ 。

\end{workedexamples}

\begin{keytakeaway}


\begin{itemize}
  \item $k$ 决定增减性： $k > 0$ 递增， $k < 0$ 递减。
  \item $b$ 决定与 $y$ 轴的交点位置。
  \item $k$ 相同则平行， $k$ 不同则相交。
\end{itemize}


\end{keytakeaway}

\begin{practice}


\begin{enumerate}
  \item $y = -\frac{1}{3}x + 2$ 中 $y$ 随 $x$ 增大是增大还是减小？
  \item 两直线 $y = 5x + 1$ 和 $y = -2x + 1$ 有什么位置关系？交点是哪里？
  \item 若一次函数图象过第一、二、三象限， $k$ 和 $b$ 的符号分别是什么？
  \item 一次函数 $y = (m-1)x + 3$ 的图象经过第一、二、三象限，求 $m$ 的取值范围。
  \item 写出一条与 $y = 2x - 1$ 平行且过点 $(0, 4)$ 的直线方程。
\end{enumerate}


\end{practice}

\section*{用函数观点看方程（组）和不等式}


\begin{explore}


方程 $2x + 1 = 0$ 的解 $x = -\frac{1}{2}$ ，在函数 $y = 2x + 1$ 的图象上对应哪个特殊的点？答案是图象与 $x$ 轴的交点 $(-\frac{1}{2}, 0)$ 。方程、不等式和函数看似不同，实际上可以用函数的图象统一理解。

\end{explore}

\begin{understand}


\textbf{函数与方程的关系}：

方程 $kx + b = 0$ 的解，就是函数 $y = kx + b$ 的图象与 $x$ 轴交点的横坐标。

\textbf{函数与不等式的关系}：

\begin{itemize}
  \item $kx + b > 0$ 的解集，就是函数 $y = kx + b$ 的图象在 $x$ 轴\textbf{上方}的部分对应的 $x$ 的范围。
  \item $kx + b < 0$ 的解集，就是图象在 $x$ 轴\textbf{下方}的部分对应的 $x$ 的范围。
\end{itemize}


\textbf{函数与方程组的关系}：

二元一次方程组的解，就是两条直线的\textbf{交点坐标}。

\end{understand}

\begin{workedexamples}


\textbf{例 1.} 利用函数图象解方程 $2x - 4 = 0$ 。

\textbf{解：} 画 $y = 2x - 4$ 的图象。与 $x$ 轴的交点为 $(2, 0)$ ，所以 $x = 2$ 。

\textbf{例 2.} 利用函数图象解不等式 $2x - 4 > 0$ 。

\textbf{解：} $y = 2x - 4$ 的图象在 $x$ 轴上方的部分对应 $x > 2$ 。所以解集为 $x > 2$ 。

\textbf{例 3.} 利用函数图象解方程组 $\begin{cases} y = x + 1 \\ y = -2x + 7 \end{cases}$ 。

\textbf{解：} 画两条直线，找交点。

由 $x + 1 = -2x + 7$ ， $3x = 6$ ， $x = 2$ ， $y = 3$ 。

两条直线的交点为 $(2, 3)$ ，即方程组的解为 $x = 2, y = 3$ 。

\textbf{例 4.} 如图，直线 $y = x - 1$ 和 $y = -\frac{1}{2}x + 2$ 相交于点 $A(2, 1)$ 。直接从图象中读出不等式 $x - 1 > -\frac{1}{2}x + 2$ 的解集。

\textbf{解：} 不等式 $x - 1 > -\frac{1}{2}x + 2$ 就是" $y = x - 1$ 在 $y = -\frac{1}{2}x + 2$ 上方"的部分。从图象可知，当 $x > 2$ 时第一条直线在第二条上方。解集为 $x > 2$ 。

\end{workedexamples}

\begin{keytakeaway}


\begin{itemize}
  \item 方程的解 = 函数图象与 $x$ 轴交点的横坐标。
  \item 不等式的解集 = 函数图象在 $x$ 轴上方（或下方）部分对应的 $x$ 范围。
  \item 方程组的解 = 两函数图象的交点坐标。
  \item 函数观点将方程、不等式、方程组统一起来。
\end{itemize}


\end{keytakeaway}

\begin{practice}


\begin{enumerate}
  \item 利用函数 $y = 3x - 6$ 的图象，求方程 $3x - 6 = 0$ 的解。
  \item 利用上题函数图象，求不等式 $3x - 6 \leq 0$ 的解集。
  \item 利用图象法解方程组 $\begin{cases} y = 2x - 1 \\ y = -x + 5 \end{cases}$ 。
  \item 一次函数 $y = kx + 3$ 的图象与 $x$ 轴的交点为 $(6, 0)$ ，求 $k$ 。
  \item 直线 $y_1 = x + 2$ 和 $y_2 = 3x - 4$ 交于一点，当 $x$ 为何值时 $y_1 > y_2$ ？
\end{enumerate}


\end{practice}



\begin{answer}


\textbf{函数的概念 练一练}

\begin{enumerate}
  \item 自变量 $r$ ，因变量 $S$ 。定义域 $r > 0$ 。
  \item 分母 $x - 1 \neq 0$ ，定义域为 $x \neq 1$ 。
  \item 不是函数。当 $x = 4$ 时， $y = 4$ 或 $y = -4$ ，不唯一。
  \item $x = -1$ ： $y = -5$ ； $x = 0$ ： $y = -2$ ； $x = 1$ ： $y = 1$ ； $x = 2$ ： $y = 4$ 。
  \item $y = 2(x - 3) + 10 = 2x + 4$ （ $x > 3$ ）。
\end{enumerate}


\textbf{函数的表示方法 练一练}

\begin{enumerate}
  \item $x = -1$ ： $y = 4$ ； $x = 0$ ： $y = 3$ ； $x = 1$ ： $y = 2$ ； $x = 2$ ： $y = 1$ ； $x = 3$ ： $y = 0$ ； $x = 4$ ： $y = -1$ 。
  \item 图象是一条直线（从左上到右下）。
  \item $y$ 值依次为 $4, 1, 0, 1, 4$ 。图象不是直线，是一条开口向上的抛物线。
  \item 解析式法：精确便于计算。列表法：查找方便。图象法：直观展示趋势。
  \item $y = 100 - 5(x - 1) = -5x + 105$ （ $1 \leq x \leq 20$ ，且 $x$ 为正整数时有意义）。或简化为 $y = -5x + 105$ 。
\end{enumerate}


\textbf{正比例函数 练一练}

\begin{enumerate}
  \item $y = -2x$ 。
  \item $k = 5 > 0$ ，图象过第一、三象限。
  \item $y = kx$ ， $-6 = 2k$ ， $k = -3$ 。 $y = -3x$ 。
  \item $-12 = 3k$ ， $k = -4$ 。
  \item $y = 3x$ 更陡，因为 $\lvert 3 \rvert > \lvert \frac{1}{2} \rvert$ 。
\end{enumerate}


\textbf{一次函数 练一练}

\begin{enumerate}
  \item 取 $(0, 4)$ 和 $(4, 0)$ 两点画直线。
  \item 与 $y$ 轴交点：令 $x = 0$ ， $y = -3$ ，交点 $(0, -3)$ 。与 $x$ 轴交点：令 $y = 0$ ， $\frac{1}{2}x = 3$ ， $x = 6$ ，交点 $(6, 0)$ 。
  \item $y = kx + b$ 。由 $(0, -1)$ 得 $b = -1$ 。由 $(2, 3)$ 得 $2k - 1 = 3$ ， $k = 2$ 。 $y = 2x - 1$ 。
  \item 经过第一、二象限说明 $b > 0$ ，经过第二象限但不经过第四象限（"下降"不够远或"上升"经过）。 $k > 0$ ， $b > 0$ 。
  \item $y = 20 - 2x$ （ $0 \leq x \leq 10$ ）。
  \item $2x + 1 = -x + 7$ ， $3x = 6$ ， $x = 2$ ， $y = 5$ 。交点 $(2, 5)$ 。
\end{enumerate}


\textbf{一次函数的图象与性质 练一练}

\begin{enumerate}
  \item $k = -\frac{1}{3} < 0$ ， $y$ 随 $x$ 增大而减小。
  \item $k$ 不同（ $5$ 和 $-2$ ），两直线相交。 $b$ 相同（都是 $1$ ），交于 $y$ 轴上的点 $(0, 1)$ 。
  \item 过第一象限说明右端上升， $k > 0$ 。过第二象限到第三象限说明与 $y$ 轴交于正半轴， $b > 0$ 。
  \item $m - 1 > 0$ （递增）且 $b = 3 > 0$ ，所以 $m > 1$ 。
  \item 平行则 $k$ 相同， $k = 2$ 。过 $(0, 4)$ 则 $b = 4$ 。直线方程 $y = 2x + 4$ 。
\end{enumerate}


\textbf{用函数观点看方程（组）和不等式 练一练}

\begin{enumerate}
  \item 令 $y = 0$ ： $3x - 6 = 0$ ， $x = 2$ 。
  \item $y \leq 0$ 对应的 $x$ 范围： $3x - 6 \leq 0$ ， $x \leq 2$ 。
  \item $2x - 1 = -x + 5$ ， $3x = 6$ ， $x = 2$ ， $y = 3$ 。解为 $x = 2, y = 3$ 。
  \item 图象过 $(6, 0)$ ： $6k + 3 = 0$ ， $k = -\frac{1}{2}$ 。
  \item $x + 2 > 3x - 4$ ， $-2x > -6$ ， $x < 3$ 。当 $x < 3$ 时 $y_1 > y_2$ 。
\end{enumerate}


\end{answer}
